\section[nosectionpage]{Ordinary \texttt{SocketServer} objects and the MQTT-Broker application}

\begin{frame}[fragile]{Ordinary \texttt{SocketServer} objects and the MQTT-Broker application}
  \begin{itemize}
    \item The MQTT broker \texttt{\textcolor{fh-ooe-red}{mqttbroker}} is an ordinary \textcolor{fh-ooe-red}{TCP server} application communicating \textcolor{fh-ooe-red}{unencrypted} via \textcolor{fh-ooe-red}{IPv4} with the MQTT broker protocol attached.
    Thus you have to use one of the \textcolor{fh-ooe-red}{template classes \texttt{SocketServer}} provided by SNode.C and provide a \texttt{\textcolor{fh-ooe-red}{SocketContextFactory}} class as \textcolor{fh-ooe-red}{template argument}.

    \begin{cppcode}[gobble=4]
      #include "net/in/stream/legacy/SocketServer.h" // stream => TCP (dgram would be UDP)
      #include "ASocketContextFactory.h"

      int main(int argc, char* argv[]) {
        ...
        net::in::stream::legacy::SocketServer<ASocketContextFactory> server(...);
        ...
      }
    \end{cppcode}
    \item The \texttt{\textcolor{fh-ooe-red}{SocketContextFactory}} is responsible for \textcolor{fh-ooe-red}{creating} a concrete \texttt{\textcolor{fh-ooe-red}{SocketContext}} object for \textcolor{fh-ooe-red}{every connected client}, \textcolor{fh-ooe-red}{representing} the \textcolor{fh-ooe-red}{application protocol} (e.g. \texttt{\textcolor{fh-ooe-red}{HTTP}}, or \texttt{\textcolor{fh-ooe-red}{Websocket}}, or \texttt{\textcolor{fh-ooe-red}{MQTT}}, or \ldots) used by the application.
  \end{itemize}
\end{frame}
